\ \\\textbf{Утверждение} Натуральное число представимо в виде суммы двух квадратов тогда и только тогда, когда любое простое число вида 4k + 3 входит в его разложение на простые множители в четной степени.
\\
\Proof 
\\
$\Rightarrow$ Пусть $n$ = $a^2$ + $b^2$. Тогда n =  $(a + i b)(a - i b) = z\overline{z}$. Разложим $z$ на простые гауссовы числа: $z = up_1...p_s * q_1...q_l$, где $p_1...p_s$ -- простые целые числа вида 4k+3, а $q_1...q_l$ -- простые гауссовы числа, имеющие ненулевую мнимую часть. Такое разложение возможно, так как все простые числа имеют вид либо 4k+3, либо 4k+1, и последнее представимо в виде произведения гауссовых чисел и u, $u \in \{\pm 1, \pm i\}$. Тогда имеем:
\begin{center}
    $\overline{z} = \overline{up_1....p_sq_1....q_l} = p_1...p_s\overline{u}\ \overline{q_1}....\overline{q_l}$ 
\end{center}
Следовательно $n = z\overline{z} = p_1^2...p_s^2u \overline{u} q_1\overline{q_1}...q_l\overline{q_l} = p_1^2...p_s^2N(q_1)...N(q_l)$ - разложение на простые множители. \\
$\Leftarrow$ $n = p_1^2...p_s^2N(q_1)...N(q_l)$. Пусть $z = p_1..p_s q_1...q_l \longrightarrow \overline{z} = p_1...p_s \overline{q_1}...\overline{q_l}$ и $n = z\overline{z} = a^2 + b^2$, если $z = a+ i b$ \EndProof
